\chapter{Complex Sentences and Clause Chaining}\label{ch:complex}

\section{Theoretical framework for complex sentences: subordination types
and the clause-combining hierarchy}
\label{sec:th-complex-overview}

The major strategies for combining clauses in Iridian reflect a unified
structural principle: \emph{head-finality at all levels of embedding}. A
dependent clause, whether it functions as an adverbial modifier, an argument,
or an attributive modifier, precedes the higher constituent whose head it is
the complement of. This property is not merely linear: it is a structural
entailment of head-final phrase structure applied recursively across clausal
levels. Table~\ref{tab:th-combining-strategies} gives the main combining
strategies and their formal properties.

\begin{table}[ht]
	\footnotesize\sffamily
	\caption{Main clause-combining strategies in Iridian.}\label{tab:th-combining-strategies}
	\medskip
	\begin{tabularx}{\textwidth}{l L L l}
		\toprule
		{\sc strategy} & {\sc formal means} & {\sc function} & {\sc ref.} \\
		\midrule
		Converb chaining &
		Non-finite converb suffixes (\ird{-u}, \ird{-ení}, \ird{-ěc},
		\ird{-ic}, \ird{-it}, \ird{-ouc}, \ird{-ema}, etc.) &
		Adverbial subordination: temporal, causal, conditional,
		concessive, purposive, contrastive &
		§\,\ref{sec:converb-chains} \\[4pt]
		Finite \ird{po}-complementation &
		Finite embedded clause + linker \ird{-e} + \ird{po} &
		Propositional complementation: saying, thinking, knowing,
		believing, etc. &
		§\,\ref{sec:complementation} \\[4pt]
		Non-assertive \ird{poby} &
		Finite embedded clause + linker \ird{-e} + \ird{poby} &
		Non-assertive finite complementation: embedded polar issues,
		embedded wh-content questions, unresolved alternatives,
		attributed non-assertive content &
		§\,\ref{sec:poby} \\[4pt]
		Purposive converb complement &
		Purposive converb + topic-coreferent subject &
		Event complementation: wanting, trying, intending, beginning &
		§\,\ref{sec:purposive} \\[4pt]
		Coordination &
		\ird{a}, \ird{še}, \ird{má}/\ird{ozná}, disjunctives &
		Symmetric clause-level coordination &
		§\,\ref{sec:coordination} \\[4pt]
		Relativization &
		Relative-converb construction (see below) &
		Nominal modification by a clausal relative &
		§\,\ref{sec:th-relativization} \\
		\bottomrule
	\end{tabularx}
\end{table}

\subsection{Head-finality in clause combining}
\label{sec:th-subord-headfinal}

The unmarked pattern in clause combining is as follows:

\begin{center}
\term{[subordinate clause] \quad [matrix-clause material] \quad [matrix verb]}
\end{center}

The subordinate clause occupies the left side of the construction; the matrix
predicate closes it. This ordering holds across all of the major subordinate
clause types: converb chains, finite complement clauses, embedded wh-clauses,
and relative clauses all precede the higher predicate they are embedded under.
The same head-final logic applies inside the subordinate clause itself,
which is internally verb-final.

\begin{figure}[ht]
	\centering
	\begin{forest}
		for tree={
			align=center,
			parent anchor=south,
			child anchor=north,
			l sep=12pt,
			s sep=16pt,
		}
		[S
			[{XP\\{\footnotesize subordinate clause}}]
			[VP
				[{XP\\{\footnotesize matrix-clause material}}]
				[{V\\{\footnotesize matrix}}]
			]
		]
	\end{forest}
	\caption{Head-final organization of subordinate constructions.}
\end{figure}

There is a practical consequence for information structure. In a converb
chain, all aspectual contrasts and alignment contrasts are borne by the
single finite verb at the end of the chain; converbs are aspectually
neutral and take nominative-accusative alignment regardless of the aspect
that the matrix verb will bear. This creates an asymmetry between
converb-internal and matrix-level information marking that is discussed in
§\,\ref{sec:converb-chains}.

\subsection{Subordination types: a typological note}
\label{sec:th-subord-types}

Following standard typological distinctions (Haspelmath
1995; Nikolaeva 2007), Iridian
clause-combining strategies may be classified along two parameters: the
\emph{morphological type of the dependent verb} (finite or non-finite) and
the \emph{semantic relationship} between dependent and matrix clause.

Non-finite (converb) subordination is the dominant strategy for adverbial
clauses. Finite subordination by means of \ird{po}/\ird{poby} is restricted
to argument clauses (complement clauses). There is no dedicated relative
clause morphology comparable to an overt relative pronoun or a relative
complementizer: relativization is expressed through a converb-based
construction in which a non-finite clause modifies the head noun (see
§\,\ref{sec:th-relativization}).

The dominance of converb-based strategies means that what in a European
language would be expressed as two (or more) finite clauses typically
appears in Iridian as a converb chain followed by a single finite verb. The
aspectual and alignment information that would be distributed across multiple
finite verbs in a Slavic language is therefore concentrated at the right edge
of the chain. This property has implications for switch-reference tracking
(§\,\ref{sec:switch-reference}) and for discourse structure
(Chapter~\ref{ch:discourse}).

\section{Converb chains: structure and alignment}\label{sec:converb-chains}

\subsection{The finite-verb-final rule}

In a converb chain, all verbs except the last are non-finite (converbs). Only the final verb is finite and carries aspect, mood, and alignment. This rule is absolute in standard written Iridian; in colloquial and dialectal speech, finite verbs may occasionally appear in non-final positions in certain coordinated structures.

\subsection{Accusative alignment within converb clauses}

Case marking within converb clauses always follows nominative-accusative alignment: the absolutive argument is the agent or sole argument, and the patient appears in the indirect (oblique) case. The ergative case is never found inside a converb clause, because converbs carry no aspect and the ergative alignment is triggered only by the perfective and retrospective aspects.

\subsection{Switch-reference and topic continuity}\label{sec:switch-reference}

By default, the subject of a converb clause is coreferent with the subject of the following clause (same-subject, SS reading). When the subjects differ (different-subject, DS), the contrast is signaled by overt lexical NPs, demonstratives, or the discourse structure itself; Iridian has no grammaticalized switch-reference clitic. Where both clauses supply full nominal subjects, disambiguation is immediate; where one subject is omitted, recoverability from discourse context determines the SS or DS reading.

\section{Sequential constructions}\label{sec:sequential}

\subsection{Forms: \ird{-u} and \ird{-ení}}

The sequential converb has two forms with subtly different temporal force. The \ird{-u} form implies a temporally bounded, completed sub-event preceding the main clause event. The \ird{-ení} form is temporally underspecified and is compatible with both sequential and simultaneous readings.

\pex
\a
\begingl
\gla Marek zdravení Janek znovime.//
\glb Marek sleep-\Seq{} Janek study-\Prog{}//
\glft \trsl{While Marek sleeps, Janek is studying.} \lingcontext{simultaneous reading with \ird{-ení}}//
\endgl
\xe

\subsection{The copula converb and the discourse particle \ird{nu}}\label{sec:seq-copula}

The marked copular verb \ird{neh-} (see §\,\ref{sec:neh} of Chapter~\ref{ch:simple}) has the sequential converb form \ird{nesu}, a suppletive form established independently of the regular aspect paradigm. In colloquial speech, \ird{nesu} is contracted to \ird{nu}. The form \ird{nu} has been grammaticalized in some dialects as a discourse particle inserted between clauses in the manner of English `and then', without strong sequential implication and without retaining the inchoative/copular force of the underlying verb.

\subsection{Interaction with applicative and causative}

\pex
\a \ird{piašt-u} / \ird{piašt-ení} \trsl{having eaten / and [then] ate\ldots}
\a \ird{piašt-éb-u} / \ird{piašt-éb-ní} \trsl{having fed someone}
\a \ird{piašt-nau} / \ird{piašt-na-ní} \trsl{having made someone eat}
\xe

\section{Temporal constructions}\label{sec:temporal}

\subsection{The simultaneous / manner converb (\ird{-ěc})}

The simultaneous converb (\ird{-ěc}) describes an action that is ongoing or co-occurring with the main clause event, and also functions as a manner converb. The suffix belongs to the palatalizing domain of front-vocalic morphology (see Section~\ref{sec:front-palatalization}), not to the special sibilant-replacement system of the antipassive. When the converb and main clause have different subjects, the sequential \ird{-ení} is preferred for purely simultaneous meaning.

This converb should not be confused with the derived manner-adverb suffix
\ird{-ě} discussed in §\,\ref{sec:stative-adv}. The converb in \ird{-ěc}
remains a non-finite verbal form used for clause-linking and event-structural
modification, whereas forms such as \ird{skorně} \trsl{quickly} are lexical
adverbs, not reduced clauses and not members of the converb paradigm.

\subsection{The six temporal converbs}

\begin{table}[ht]
	\footnotesize\sffamily
	\caption{Temporal converb suffixes.}\label{tab:tempconverbs}
	\medskip
	\begin{tabularx}{0.9\textwidth}{L l L}
		\toprule
		{\sc suffix} & {\sc meaning} & {\sc function} \\
		\midrule
		\ird{-ezak}    & until              & main clause holds until converb-clause event \\
		\ird{-eškady}  & when, around when  & loose temporal co-occurrence \\
		\ird{-eeby}    & since              & from the time of the converb-clause event \\
		\ird{-ešhoume} & as soon as         & immediately upon the converb-clause event \\
		\ird{-eebymy}  & after              & after the converb-clause event \\
		\ird{-esny}    & before             & before the converb-clause event \\
		\bottomrule
	\end{tabularx}
\end{table}

\section{Causal constructions}\label{sec:causal}

% ============================================================
% STATUS: STUB — causal converb suffix not yet documented.
% The 2026 notes list converb types 1, 2, 4, 5 and skip 3 (causal).
% This section is a placeholder until the suffix is confirmed.
% ============================================================

A causal converb introduces the reason or cause for the main clause event (`because X, Y'). The suffix is not yet confirmed in the current revision and will be documented here once established.

\section{Concessive constructions}\label{sec:concessive}

The concessive converb (\ird{-ouc}) introduces a conceded or contrasted clause (`although X, Y'). It licenses the irrealis mood in the main clause when the conceded content is contrary to expectations.

\section{Conditional constructions}\label{sec:conditional}

\subsection{The four conditional converbs}

\begin{table}[ht]
	\footnotesize\sffamily
	\caption{Conditional converb suffixes.}\label{tab:conditional}
	\medskip
	\begin{tabularx}{0.8\textwidth}{L l L}
		\toprule
		{\sc suffix} & {\sc meaning} & {\sc notes} \\
		\midrule
		\ird{-ic}   & if (possibility)       & does not trigger palatalization \\
		\ird{-emy}  & if not (possibility)   & \\
		\ird{-ezmy} & if/when (expectation)  & \\
		\ird{-ena}  & if/when not (expected) & \\
		\bottomrule
	\end{tabularx}
\end{table}

\subsection{No palatalization with \ird{-ic}}\label{sec:cond-nopala}

Despite beginning with the front vowel \ird{-i-}, the conditional converb \ird{-ic} does not trigger palatalization of the thematic consonant. This is a historical accident: the suffix was grammaticalized during a period when this \ird{-i-} lacked palatalizing force. The conditional is therefore distinguishable from the irrealis (\ird{-ěv}), the simultaneous converb (\ird{-ěc}), and the conjunctive linker \ird{-e}, all of which belong to the productive palatalizing domain.

\section{Purposive constructions}\label{sec:purposive}

The purposive converb (\ird{-it}) expresses purpose (`in order to X') and serves as the complement of modal and desiderative verbs. Palatalization is irregular, applying mainly to Class~I stems.

\subsection{The intentional future reanalysis}

In colloquial speech, the purposive converb has been reanalyzed as a future/intentional marker when no governing verb is present. The first-person subject pronoun \ird{dá} is typically retained overtly in this construction, as the purposive suffix alone does not encode person:

\ex
\begingl
\gla Dá trave piašcit.//
\glb \textsc{1sg} bread.\Ind{} eat-\Purp{}//
\glft \trsl{I'm gonna eat bread.} \lingcontext{colloquial intentional future}//
\endgl
\xe

\section{Contrastive constructions}\label{sec:contrastive}

The contrastive converb (\ird{-ema}) introduces a contrasting clause (`X, but/whereas Y'), covering simple contrast, adversative contrast, and parallel contrast.

\section{Coordination}\label{sec:coordination}

% ============================================================
% STATUS: STUB — systematic treatment needed.
% Key conjunctions: a (additive), še (additive / complementizer),
% má / ozná (adversative), disjunctives.
% ============================================================

Coordination is expressed by \ird{a} (additive), \ird{še} (additive, also complementizer), \ird{má} / \ird{ozná} (adversative), and a set of disjunctives.

\section{Relativization}
\label{sec:th-relativization}

\subsection{Structure of the relative construction}
\label{sec:th-rel-structure}

Iridian lacks a dedicated relative pronoun or overt relative complementizer.
Relative clauses are formed using a converb-based strategy: the verb of the
relative clause bears a relativizing non-finite form and the entire relative
clause precedes the head noun it modifies, consistent with the general
head-final principle. The head noun appears after the relative clause in the
surface string.

\begin{center}
\term{[relative clause + \textsc{rel.form}] \quad [head noun]}
\end{center}

The structure is thus the mirror image of postnominal relative constructions
in European languages. Where English says \textit{the man who bought the
book}, Iridian places the relative clause before its head: the element
corresponding to \textit{who bought the book} precedes the element
corresponding to \textit{the man}.

\subsection{Head-finality and the placement of relative clauses}
\label{sec:th-rel-headfinal}

Because the relative clause is an attributive modifier, it modifies the
head noun from which it is separated only by the nominal head. This prenominal
position is the expected one under head-finality: just as adjectival modifiers
and demonstratives precede the noun they modify in other nominal constructions,
so the clausal relative modifier precedes its head. No extraposition or
rightward movement of relative clauses occurs in standard Iridian.

The head noun itself assigns case to its argument position in the matrix clause,
independently of whatever role the coreferent argument plays inside the relative
clause. Case-mismatch between the head noun's matrix role and its relative-clause
role is resolved by the relativization construction itself rather than by any
case variation on the head noun.

\subsection{The gap, resumption, and argument-relativization hierarchy}
\label{sec:th-rel-gap}

% ============================================================
% STATUS: STUB — relativization morphology not yet fully
% documented in the 2026 rewrite. The following describes the
% expected properties; update with attested data once the
% relativization section of ch.\ 04 is written.
% ============================================================

The relativized argument position inside the relative clause is typically
represented by a gap (silent trace). The accessibility hierarchy of
relativization positions is expected to follow the standard cross-linguistic
pattern (Keenan \& Comrie 1977), with subjects and direct objects the most
freely relativizable positions. Indirect objects and oblique arguments may
employ a resumptive pronoun strategy in contexts where a gap would yield
structural ambiguity; this expectation should be verified against corpus
data.

Relativization of the locative and other oblique positions proceeds with the
postpositional phrase intact and the oblique noun gapped or resumed within the
relative clause. The head noun of such a construction is the entire gapped-PP
antecedent; the postposition does not strand to the right of the head.


\subsection{}

Iridian possesses a special non-finite suffix -eň, used with nominal predicates
and with a restricted class of stative predicates to form state-frame adverbial
clauses. The construction expresses meanings such as ‘while X’, ‘when in the
state of X’, and ‘as X’. It is functionally parallel to the ordinary converb
system in that it forms a preverbal adverbial dependent clause, but it is not
part of the regular converb inventory and does not replace an ordinary Asp.Align
morpheme in the verbal paradigm.

Semantically, -eň backgrounds a classification or state as the temporal or
circumstantial frame for the matrix event. The clause does not assert entry into
the state, progression toward it, or completion of a transition. In this it
contrasts sharply with neh-, whose function is explicitly transitional and
inchoative: neh- introduces an event of becoming and projects a full verbal
spine under AspP. A form such as \ird{byl nesu} would therefore mean something
like ‘while becoming a child’ or ‘while entering child-status’, not the intended
simple background reading ‘while a child’. The state-frame converb -eň is the
dedicated exponent of the latter meaning.

Syntactically, a -eň clause behaves like other Iridian adverbial subordinate clauses. It precedes the matrix clause in accordance with general head-finality, and its interpretation is normally same-subject unless context or overt nominal material indicates otherwise. Because the construction is not an ordinary finite clause, it bears no independent aspect or alignment morphology; the matrix finite verb remains the sole carrier of aspect, mood, and alignment in the chain, in accordance with the general finite-verb-final rule for Iridian subordination.


The suffix is primarily used with class-membership nouns and age/status predicates:
byl-eň ‘while a child, as a child’
kral-eň ‘while king, as king’
vdov-eň ‘while a widow’

The phasal particle s may occur with the state-frame converb -eň. In such cases it scopes over the predicative state expressed by the converb and marks that state as still holding at the reference time of the subordinate clause. Thus byleň means ‘while a child’, whereas s byleň means ‘while still a child’. Any implication that the state no longer holds at speech time arises pragmatically from discourse perspective, not from the particle itself.


\section{Clausal complementation}\label{sec:complementation}

\subsection{The propositional/event-complement distinction}
\label{sec:th-compl-distinction}

Iridian distinguishes two major types of clausal complementation according to
the semantic type of the embedded content: \term{propositional complementation}
and \term{event complementation}. This distinction, which parallels the
broad cross-linguistic division between that-clause and to-infinitive
complements in English, is encoded formally by different constructions in
Iridian.

Propositional complementation is expressed by the finite \ird{po}-clause:
the embedded clause remains fully finite, retaining all verbal categories
of an independent clause, and is subordinated by the linker \ird{-e}
followed by \ird{po}. Event complementation is expressed by the purposive
converb: the embedded verb is non-finite and the construction is semantically
of the controlled or prospective type.

\begin{description}[leftmargin=2em, labelwidth=!, labelindent=0em]
\item[\upshape Propositional complement (\ird{po}):] used with predicates of
	saying, reporting, thinking, knowing, believing, remembering, inferring,
	perceiving a proposition, and evaluating a proposition. The complement
	presents a state of affairs as held to be true (or false, or probable,
	etc.).\ by the matrix subject.
\item[\upshape Event complement (purposive converb):] used with predicates of
	wanting, trying, intending, beginning, going (purposively), and related
	predicates of desire, attempt, inception, and motion. The complement
	presents an event as projected, volitional, incipient, or
	purpose-oriented rather than as a proposition held to be true.
\end{description}

\noindent The distinction is not merely semantic: it has formal consequences.
Propositional complements are always finite; event complements are always
non-finite. In languages where the distinction collapses (as it does for many
predicates in colloquial registers), Iridian preserves it morphologically.

\subsection{The finite \ird{po}-construction}

The ordinary means of embedding a full proposition is the finite
\ird{po}-clause. Its form is:

\begin{center}
\term{[finite embedded clause + linker \ird{-e}] \ird{po}}
\end{center}

The embedded clause remains fully finite. It preserves the same verbal
categories as an independent clause, including aspect, the alignment
distinctions associated with aspect, voice, and, where required by the sense,
mood. The construction is therefore not a reduced verbal form, but a true
finite clause made subordinate by the linker \ird{-e} and the complement
marker \ird{po}.

An independent clause such as the following:

\ex
\begingl
\gla Lobek Marku piaštnik.//
\glb apple.\Dir{} Marek.\Trs{} eat-\Pf{}//
\glft \trsl{Marek ate the apple.}//
\endgl
\xe

may be embedded without loss of finiteness:

\ex
\begingl
\gla Lobek Marku piaštnice po Janku dalnik.//
\glb apple.\Dir{} Marek.\Trs{} eat-\Pf{}-\Lnk{} \Compl{} Janek.\Trs{} say-\Pf{}//
\glft \trsl{Janek said that Marek ate the apple.}//
\endgl
\xe

The same construction is used after predicates of saying, reporting,
thinking, knowing, believing, remembering, inferring, perceiving a
proposition, and evaluating a proposition. In all such cases the complement
is construed as a full proposition, not merely as a projected event.

The \ird{po}-clause may also appear without an overt matrix predicate. Such
forms were formerly described as \textit{bare quotatives}; synchronically it
is better to treat them as matrixless uses of the same finite complement
construction. In such cases the proposition is presented not as a simple
speaker-assertion, but as attributed content. According to context, this may
yield a reportive, hearsay, echoic, reflective, distancing, ironic, or
reservation-marking effect.

\ex
\begingl
\gla Lobek Marku piaštnice po.//
\glb apple.\Dir{} Marek.\Trs{} eat-\Pf{}-\Lnk{} \Compl{}//
\glft \trsl{Marek ate the apple, I hear.}//
\endgl
\xe

The precise omitted frame need not be supplied in every instance. One may
understand such clauses loosely as equivalent to \textit{it is said},
\textit{I hear}, \textit{one gathers}, \textit{so they say}, or the like;
but grammatically the essential fact is that the proposition is cast as
attributed rather than directly asserted.

\subsection{Prospective and controlled event complements}

The finite \ird{po}-construction must be kept distinct from the complement
used for intended, attempted, controlled, or projected events. This second
construction belongs to the domain of the purposive converb. It is employed
with predicates such as \textit{want}, \textit{try}, \textit{intend},
\textit{go in order to}, \textit{begin}, and similar verbs whose complement
denotes not a proposition as such, but an event regarded as prospective,
volitional, incipient, or purposive.

In this use the complement is non-finite. It does not present an independent
state of affairs as true, reported, believed, or remembered, but rather an
event toward which the matrix predicate is directed. The distinction is
therefore one between propositional complementation and event
complementation. The former is the province of finite \ird{po}-clauses; the
latter, of the purposive converb.

\subsection{The non-assertive complement \ird{poby}}\label{sec:poby}

\subsubsection{The \ird{po} / \ird{poby} opposition}

The fundamental contrast between \ird{po} and \ird{poby} is one of assertive
versus non-assertive complementation. \ird{Po} is the assertive
complementizer: it introduces a finite embedded clause presented as
propositional content --- something said, believed, remembered, known, or
otherwise treated as an assertable proposition. \ird{Poby} is the
non-assertive complementizer: it introduces a finite embedded clause presented
as an issue, question, unresolved proposition, attributed but unendorsed
content, or other non-assertively framed material. The form in both cases is:

\begin{center}
\term{[finite embedded clause + linker \ird{-e}] \ird{po} / \ird{poby}}
\end{center}

The embedded clause remains fully finite in both constructions. Both \ird{po}
and \ird{poby} are therefore best analysed as heads of CP selecting a clause
that may itself project \textsc{PolP}; the complementizer contrast is a
CP-level phenomenon (§\,\ref{sec:vm-tension-linker} in
Chapter~\ref{ch:verbs}).

\subsubsection{Three complement types under \ird{poby}}

The non-assertive framing of \ird{poby} covers three complement types, which
are not three unrelated constructions but three values of the same broader
category.

\paragraph{Type~1: Embedded polar issue (\textit{whether}-complement).}
The first and most central use embeds an unresolved polar question: whether a
proposition holds. This is the closest Iridian equivalent to English
\textit{whether}. The embedded clause is a full finite clause without a
wh-operator. After verbs of asking, wondering, doubting, or deciding,
\ird{poby} marks the embedded content as an open issue rather than an
asserted fact.

\pex
\a
\begingl
\gla Lobek Marku piaštnice poby Janku houstnik.//
\glb apple.\Dir{} Marek.\Trs{} eat-\Pf{}-\Lnk{} \Compl{} Janek.\Trs{} ask-\Pf{}//
\glft \trsl{Janek asked whether Marek ate the apple.}//
\endgl
\a
\begingl
\gla Marek ščene poby Janek rodžounik.//
\glb Marek.\Dir{} arrive-\Pf{}-\Lnk{} \Compl{} Janek.\Dir{} wonder-\Cont{}//
\glft \trsl{Janek is wondering whether Marek arrived.}//
\endgl
\xe

\paragraph{Type~2: Embedded wh-content question.}
The second type embeds a content question: who came, what he bought, where
they went. Here the embedded clause contains a wh-operator from the inventory
in Table~\ref{tab:nm-interrogatives} (Chapter~\ref{ch:nouns}), and \ird{poby}
marks the clause as non-assertively embedded. The embedded wh-phrase
ordinarily remains in situ in the neutral case; optional fronting within the
embedded clause's own left periphery signals marked focus, consistent with the
general principle established in §\,\ref{sec:th-wh-interrogation}.

\pex
\a
\begingl
\gla Jede lobek Marku piaštnice poby Janek rodžounik.//
\glb who apple.\Dir{} Marek.\Trs{} eat-\Pf{}-\Lnk{} \Compl{} Janek.\Dir{} wonder-\Cont{}//
\glft \trsl{Janek is wondering who Marek gave the apple to.} \lingcontext{wh in situ}//
\endgl
\a
\begingl
\gla Marek jena mulše poby Janek houstnik.//
\glb Marek.\Dir{} where live-\Cont{}-\Lnk{} \Compl{} Janek.\Trs{} ask-\Pf{}//
\glft \trsl{Janek asked where Marek lives.} \lingcontext{wh in situ}//
\endgl
\xe

\paragraph{Type~3: Attributed non-assertive content.}
The third type embeds finite clauses under verbs of saying, reporting,
suggesting, discussing, or similar predicates when the matrix speaker does not
present the embedded content as a straightforward asserted proposition but
rather as attributed, unendorsed, reported, dubitatively framed, or otherwise
held at arm's length. This is the non-assertive quotation function of
\ird{poby}: contrast the following minimal pair with \ird{po}:

\pex
\a
\begingl
\gla Marek ščenice po Janek vlašnik.//
\glb Marek.\Dir{} arrive-\Pf{}-\Lnk{} \textsc{assert.compl} Janek.\Dir{} report-\Pf{}//
\glft \trsl{Janek reported that Marek arrived.} \lingcontext{\ird{po}: speaker endorses as propositional content}//
\endgl
\a
\begingl
\gla Marek ščenice poby Janek vlašnik.//
\glb Marek.\Dir{} arrive-\Pf{}-\Lnk{} \textsc{nonassert.compl} Janek.\Dir{} report-\Pf{}//
\glft \trsl{Janek reportedly said that Marek arrived.} \lingcontext{\ird{poby}: speaker attributes but does not endorse}//
\endgl
\xe

\noindent This use is especially natural under predicates of saying,
discussing, or suggesting where the matrix speaker attributes the content to
another source. It overlaps in use with the particle \ird{či} (hearsay,
§\,\ref{sec:class2-particles}), but the two are structurally distinct: \ird{či}
is a matrix-level stance marker on the main clause, while \ird{poby} is a
subordinating complementizer embedding a finite clause.

\subsubsection{Internal syntax of \ird{poby}-clauses}

In \ird{poby}-complements, the embedded clause projects \textsc{PolP}, but
neutral embedded content ordinarily uses null Pol\textsuperscript{0}. The
complementizer \ird{poby} already marks the clause as non-assertive and
issue-like; an embedded plain polar issue therefore does not require an
additional overt \ird{by} internally. Overt internal \ird{by} is reserved for
marked, echoic, or explicitly interrogative embedded readings, where the
speaker wishes to foreground the interrogative status of the embedded content
rather than presenting it as a background issue.

The architecture for the three types is thus as follows. An embedded polar
issue has a fully finite embedded clause with null Pol\textsuperscript{0}: the
non-assertive status is contributed entirely by \ird{poby}. An embedded
wh-content clause contains the wh-operator plus null Pol\textsuperscript{0}.
A marked embedded echo question may contain overt internal \ird{by} alongside
the wh-operator, in a double-marking configuration that asserts the explicitly
interrogative, insisted, or surprised character of the embedding.

\noindent This construction is found after predicates such as \textit{ask},
\textit{wonder}, \textit{inquire}, \textit{investigate}, \textit{check},
\textit{doubt}, \textit{decide}, \textit{know whether}, \textit{report
(non-assertively)}, \textit{suggest}, \textit{discuss}, and others that
take non-assertive embedded content. Predicates of the \textit{say that},
\textit{believe that}, \textit{remember that}, or \textit{know that} type
require the assertive \ird{po}-construction.

\subsection{Indirect evidentiality and attributed content}

Finite clauses in Iridian are, in themselves, evidentially unmarked. The
ordinary declarative does not grammatically specify whether the speaker knows
the matter directly or indirectly, though in practice it is often compatible
with direct knowledge. Where the source of knowledge is explicitly indirect,
Iridian employs a special indirect evidential. This marks that the
proposition rests on report, inference, visible result, or post hoc
discovery rather than on firsthand evidence. In suitable contexts a mirative
colour may arise from this meaning, especially where the information is
sudden or contrary to expectation; but the grammatical value remains that of
indirect evidence.

This category must not be confused with the matrixless \ird{po}-construction.
The indirect evidential concerns the source of knowledge. A matrixless
\ird{po}-clause concerns the mode of presentation in discourse. The two
frequently approach one another in use, since hearsay lies near the border of
both; but they are not the same thing.

Morphologically the indirect evidential belongs to the finite verbal complex.
It stands before the complement linker \ird{-e}, occupying the outermost
position in the finite inflectional spine. The order is therefore as follows:

\begin{center}
\textsc{stem} + \textsc{asp.align} + (\textsc{mood}) + (\textsc{ind.evid}) + \ird{-e} + \ird{po}/\ird{poby}
\end{center}

If the indirect evidential has the form \ird{-ejte}, then an embedded finite
clause will show:

\begin{center}
\term{[finite clause-\ird{ejte}-\ird{-e}] \ird{po}}
\end{center}

Compare the simple finite clause:

\ex
\begingl
\gla Lobek Marku piaštnicejte.//
\glb apple.\Dir{} Marek.\Trs{} eat-\Pf{}-\textsc{ind.evid}//
\glft \trsl{Marek ate the apple, apparently.}//
\endgl
\xe

When the speaker chiefly wishes to indicate that the information is
second-hand, inferred, or otherwise non-firsthand, the indirect evidential is
the ordinary means. When the speaker chiefly wishes to cast the proposition
as attributed discourse, with reportive, distancing, sceptical, ironic, or
echoic effect, the matrixless \ird{po}-construction is the ordinary means.
The two may co-occur, but such co-occurrence is marked, and yields a layered
interpretation: a proposition already marked as indirectly known is then
further presented as attributed content. This is suitable for rumour,
mediated narration, gossip, and similar discourse types, but should not be
treated as the unmarked way of expressing hearsay.

The two categories are therefore to be kept distinct. The indirect
evidential is a grammatical marker of indirect source of knowledge. The
matrixless \ird{po}-clause is a construction of attributed content, whose
pragmatic effects often include distance, reserve, scepticism, irony, or
reduced commitment, but whose core meaning is not evidential in itself.

\section{Fixed expressions with converbs}\label{sec:fixed-converbs}

% ============================================================
% Thanks, apologies, congratulations, greetings that are historically
% grammaticalized converb forms — to be documented.
% ============================================================

\section{Integrating the adverbial particle analysis: a roadmap}
\label{sec:th-particle-integration}

The adverbial particle analysis of §\,\ref{sec:adv-particles}
(Chapter~\ref{ch:function-words}) provides the primary empirical and
theoretical foundation for the left-peripheral architecture described in
this chapter. The integration works at two levels.

\emph{Empirically}, the particle ordering, scope interactions, and
fragment-licensing diagnostics establish the functional hierarchy
summarized in Table~\ref{tab:th-clause-spine} and
Figure~\ref{fig:th-left-periphery}. Any account of interrogation,
evidentiality, or discourse structure in the simple or complex sentence
chapters should treat the particle hierarchy as the structural baseline.

\emph{Theoretically}, the architecture allows a unified analysis of several
phenomena that would otherwise require separate stipulations. The licensing
of wh-operators --- whether in situ or fronted to Spec,FocP --- as
interrogative variables follows from the general principle in
§\,\ref{sec:th-multiple-wh}: a wh-expression within the scope domain of the
clause-level interrogative operator in PolP is interpreted as an interrogative
variable regardless of its surface position. Overt fronting signals marked
focus or discourse prominence; it is not required for interrogative licensing.
The surface-initial appearance of peripheral particles in topic-drop clauses
follows from null-topic rather than any separate particle-fronting rule.
The scope separation between \ird{by} (Pol\textsuperscript{0}) and the
comment-domain operators follows from their distinct projections in the clause
spine: \ird{by} occupies PolP, which is above the comment-domain
(§\,\ref{sec:by-pol-head}). And the interaction between the indirect evidential
and the matrixless \ird{po}-clause follows from their being in distinct
grammatical components (morphological verbal complex vs.\ clause-peripheral
particle).

The following cross-references indicate where the theoretical discussion in
this chapter is empirically supported or extended:

\begin{description}[leftmargin=2em, labelwidth=!, labelindent=0em]
\item[\upshape Left-periphery structure (§\,\ref{sec:th-left-periphery}):] see
	§\,\ref{sec:adv-particles} for the full particle inventory, ordering tests,
	and scope diagnostics that motivate Figures~\ref{fig:th-clause-fields}
	and~\ref{fig:th-left-periphery}.
\item[\upshape Polar questions (§\,\ref{sec:th-polar-questions}):] see
	§\,\ref{sec:by-operator} for \ird{by} as clause-typing operator and the
	predicate-focus variant; §\,\ref{sec:by-kamo} for scope independence from
	comment-domain operators.
\item[\upshape \textit{Wh}-interrogation (§\,\ref{sec:th-wh-interrogation}):] see
	§\,\ref{sec:wh-questions} for single- and multiple-wh patterns, the
	in-situ licensing principle, and the typology of question types.
\item[\upshape Evidential particles (§\,\ref{sec:th-evidential}):] see
	§\,\ref{sec:evidential-ops} (\ird{hlavdy}, \ird{ktady}, \ird{oče},
	\ird{či}) and §\,\ref{sec:complementation} for the interface between
	evidential morphology and attributed discourse.
\item[\upshape Complementation (§\,\ref{sec:th-complementation}):] see
	§\,\ref{sec:complementation} for the full treatment of \ird{po}/\ird{poby},
	the purposive converb complement, and the matrixless \ird{po}-clause.
\item[\upshape Head-finality and subordination (§\,\ref{sec:th-subord-headfinal}):] see
	§\,\ref{sec:converb-chains} for the syntax and information structure of
	converb chains, and §\,\ref{sec:sequential}--§\,\ref{sec:purposive} for
	individual subordinate types.
\end{description}
